\clearpage
\VerdierPhysicalPageAnchor{0067}%
\VerdierSemanticAnchor{ch1-2-1-5}
\noindent\textbf{2.1.5.}\quad
态射 $f^\sigma\colon Y^\sigma\to Y'^\sigma$ 称为态射 $(f^\sigma)$ 的
\emph{$\sigma$ 次分量}。设 $f\colon Y\to Y'$ 为 $n$-重复形的态射，
$f$ 的 $\sigma$ 次分量记作 $f^\sigma$。
\VerdierSemanticAnchor{ch1-2-1-5-composition}
$n$-重复形的态射以显然的方式逐分量复合。
\VerdierSemanticAnchor{ch1-2-1-5-category}
$C$ 中对象的 $n$-重复形组成一个范畴，记作 $\operatorname{comp}^{n}(C)$。
当 $n=1$ 时，我们采用简化记号
$\operatorname{comp}^{1}(C)=\operatorname{comp}(C)$。
$\operatorname{comp}(C)$ 的对象称为 $C$ 中的\emph{单复形}；在不会混淆时，
也简称为 $C$ 中的\emph{复形}。各范畴 $\operatorname{comp}^{n}(C)$ 都是
加法范畴。

\VerdierSemanticAnchor{ch1-2-1-6}
\noindent\textbf{2.1.6.}\quad
设 $Y$ 为 $C$ 中的一个 $n$-重复形。我们有时以 $Y_{-\sigma}$ 表示 $Y$ 的
$\sigma$ 次分量，并以 $d^Y_{-\sigma,i}$ 表示 $Y$ 的 $i$ 方向上的
$\sigma$ 次微分。
\VerdierSemanticAnchor{ch1-2-1-6-positive-finite}
当 $Y$ 是单复形，且其正次数分量除有限多个外都为零时，尤其使用这种记号。
同样，设 $f\colon Y\to Y'$ 为 $n$-重复形的态射；我们有时以
$f_{-\sigma}$ 表示态射 $f$ 的 $\sigma$ 次分量。

\VerdierSemanticAnchor{ch1-2-1-7}
\noindent\textbf{2.1.7.}\quad
设 $\varphi\colon[n]\to[m]$ 为一个映射。
\VerdierSemanticAnchor{ch1-2-1-7-summable}
若对每个 $\tau\in\mathbb Z[m]$，满足 $\varphi(\sigma)=\tau$ 且
$Y^\sigma\neq0$ 的 $\sigma\in\mathbb Z[n]$ 只有有限多个，则称 $C$ 中的
$n$-重复形 $Y$ 为\emph{$\varphi$-可求和的}。当 $\varphi$ 是唯一的映射
$[n]\to[1]$ 时，简称 $Y$ 为\emph{可求和的}。
\VerdierSemanticAnchor{ch1-2-1-7-composition}
设 $[n]\xrightarrow{\varphi}[m]\xrightarrow{\psi}[p]$ 为两个映射。
每个 $\psi\varphi$-可求和的 $n$-重复形 $Y$ 也都是 $\varphi$-可求和的。
特别地，每个可求和的 $n$-重复形，对任意映射
$[n]\xrightarrow{\varphi}[m]$ 都是 $\varphi$-可求和的。
\VerdierSemanticAnchor{ch1-2-1-7-injective}
当 $\varphi$ 为单射时，所有复形都是 $\varphi$-可求和的。
\VerdierSemanticAnchor{ch1-2-1-7-category}
\emph{$\varphi$-可求和 $n$-重复形的范畴}，是
$\operatorname{comp}^{n}(C)$ 中由 $\varphi$-可求和复形所确定的全子范畴。

\VerdierSemanticAnchor{ch1-2-2}
\subsection{关联单复形}

\VerdierSemanticAnchor{ch1-2-2-1}
\noindent\textbf{2.2.1.}\quad
对每个 $\sigma\in\mathbb Z[n]$，写作
$\sigma=\sum_{j\in[n]}l_je_j$，$l_j\in\mathbb Z$，并对每个
$i\in[n]$，令

\VerdierSemanticAnchor{ch1-2-2-1-eq1}
\begin{equation}
  \varepsilon(\sigma,i)=(-1)^{\left(\sum_{j<i}l_j\right)}.
  \tag{2.2.1.1}
\end{equation}

\VerdierSemanticAnchor{ch1-2-2-1-sign-map}
公式 (2.2.1.1) 定义了从 $\mathbb Z[n]\times[n]$ 到符号群
$\{-1,1\}$ 的映射
\begin{equation*}
  (\sigma,i)\longmapsto\varepsilon(\sigma,i).
\end{equation*}
立即可验证，对每个 $i\in[n]$，映射
$\sigma\longmapsto\varepsilon(\sigma,i)$ 是群同态，并且
